\subsection{Positive mass and bigness}
\label{sec:positive-mass-bigness}

The preceding equations were solved in big classes. They also detect bigness at the pseudo-effective boundary: adding $\varepsilon\omega$ first gives a big class, and an estimate uniform in $\varepsilon$ produces a strictly positive current in the original class. We begin with the dual criterion needed for that last step. Throughout this subsection $X$ is connected, compact and K\"ahler, with $n\geq1$.

\begin{lemma}[Lamari's criterion]\label{lem:kahler-lamari-criterion}
Let $\theta$ be a smooth closed real $(1,1)$-form. There exists $u\in\PSH(X,\theta)$ if and only if
\begin{equation}\label{eq:lamari-positive-tests}
 \int_X\theta\wedge\tau\geq0
\end{equation}
for every smooth strictly positive real $(n-1,n-1)$-form $\tau$ satisfying $\ddc\tau=0$. Consequently, the class of $\theta$ is big if and only if there exists $\delta>0$ such that, for all these tests,
\begin{equation}\label{eq:lamari-strict-tests}
 \int_X\theta\wedge\tau\geq\delta\int_X\omega\wedge\tau.
\end{equation}
\end{lemma}
\begin{proof}
Necessity is distributional integration by parts against $\tau$. For sufficiency let $F$ be the real vector space of smooth real $(n-1,n-1)$-forms and $S=\ker\ddc\subset F$. Put $e=\omega^{n-1}$, with $e=1$ for $n=1$, and $L(\tau)=\int_X\theta\wedge\tau$ on $S$. Adding $t e$ and letting $t\downarrow0$ shows that $L$ is nonnegative on the semipositive elements of $S$ too. Every $\eta\in F$ satisfies $-C e\leq\eta\leq C e$ for some $C$, by compactness and pointwise positivity. Define
\[
 p(\eta)=\inf\{L(\tau):\tau\in S,\ \tau\geq\eta\}.
\]
This is finite: $C e$ is admissible, whereas $\tau\geq\eta\geq-Ce$ implies $L(\tau)\geq-CL(e)$. It is sublinear by addition and positive scaling of admissible forms. If $\eta\in S$, then $\eta$ itself is admissible and every other admissible $\tau$ has $L(\tau-\eta)\geq0$, proving $p|_S=L$.

Here is the linear extension argument. Suppose $\ell\leq p$ is linear on a subspace $M\supset S$ and extends $L$. For $x\notin M$ choose a real number $a$ between
\[
 \sup_{y\in M}\bigl(\ell(y)-p(y-x)\bigr)
 \quad\text{and}\quad
 \inf_{z\in M}\bigl(p(z+x)-\ell(z)\bigr).
\]
The left endpoint is at most the right one, since
$\ell(y+z)\leq p(y+z)\leq p(y-x)+p(z+x)$; both endpoints are finite, by taking $y=z=0$ and using this same inequality. Positive scaling, separately for positive and negative coefficients of $x$, shows that $\ell(y+t x)=\ell(y)+t a$ is an extension dominated by $p$. The union along a chain of such extensions is still an extension dominated by $p$. The maximal-chain principle therefore gives an extension to all of $F$: a proper domain could be enlarged by the step just proved.

For $\eta\geq0$, the test $\tau=0$ gives $p(-\eta)\leq0$, so $\ell(\eta)\geq0$. Moreover $-Ce\leq\eta\leq Ce$ gives $|\ell(\eta)|\leq CL(e)$. Choosing $C$ bounded by a fixed multiple of the uniform coefficient norm proves continuity. Thus $\ell$ is a positive $(1,1)$-current $T$. The current $R=T-\theta$ annihilates $S$.

We verify internally that $R=\ddc f$. The forms $\partial\xi$ and $\bar\partial\xi$ of complementary bidegree are $\ddc$-closed; complexifying the annihilation identity shows $\mathrm dR=0$. Also $\int_XR\wedge\omega^{n-1}=0$, so $\Lambda R$ has mean zero. By \cref{lem:compact-elliptic-foundations,lem:kahler-trace-star}, solve $Lf=-4\pi\Lambda R$ distributionally. Then $R'=R-\ddc f$ is closed with zero trace, hence harmonic and smooth by those same lemmas. Its star $*R'$ is a smooth closed $(n-1,n-1)$-form, and therefore belongs to $S$. Stokes gives
\[
 \|R'\|_{L^2}^2
   =\int_X(R-\ddc f)\wedge *R'=0.
\]
Thus $R'=0$. In dimension one the zero trace already gives $R'=0$. Since $T=\theta+\ddc f\geq0$, local positive-current potentials and Weyl's lemma give the canonical quasi-psh representative of $f$, exactly as in the proof of \cref{thm:global-ddc-lemma}. This proves sufficiency. Apply it to $\theta-\delta\omega$ to obtain the strict criterion.
\end{proof}
\index{Lamari's criterion}\index{positive current!dual criterion}

The next estimate uses ordinary trace mass in its first factor. That trace may charge the analytic pole set, so it must not be replaced by its non-pluripolar restriction.

\Needspace{14\baselineskip}
\begin{lemma}\label{lem:prescribed-density-mixed-trace}
Let $S$ be a closed positive $(1,1)$-current whose local potentials are locally bounded on $X\setminus A$, where $A$ is a proper analytic subset. Let $\tau$ be a smooth strictly positive real $(n-1,n-1)$-form, and put $\nu=\omega\wedge\tau$. If
\[
 \langle S^n\rangle=m\nu,\qquad m\geq0,
\]
then
\begin{equation}\label{eq:prescribed-density-mixed-trace}
 \left(\int_X S\wedge\tau\right)
 \left(\int_X\langle S^{n-1}\rangle\wedge\omega\right)
               \geq\frac mn\left(\int_X\nu\right)^2.
\end{equation}
For $n=1$ the second factor is $\int_X\omega$.
\end{lemma}
\begin{proof}
For a smooth positive definite form $S$, use $\omega$ as the unit matrix at a point and diagonalize $S$, with eigenvalues $\lambda_j>0$. Let $t_j>0$ be the diagonal coefficients of the positive dual matrix of $\tau$. Then
\[
 \frac{S\wedge\tau}{\omega\wedge\tau}
       =\frac{\sum_j\lambda_jt_j}{\sum_jt_j},\qquad
 \frac{S^{n-1}\wedge\omega}{S^n}
       =\frac1n\sum_j\lambda_j^{-1}.
\]
Since $\lambda_j\sum_k\lambda_k^{-1}\geq1$, their product is at least $1/n$. Adding $t\omega$ and letting $t\downarrow0$ includes semipositive $S$. Thus, as densities with respect to any fixed smooth volume form,
\begin{equation}\label{eq:pointwise-mixed-trace}
 (S\wedge\tau)(S^{n-1}\wedge\omega)
                     \geq\frac1n S^n\nu.
\end{equation}

To pass to bounded potentials, fix a coordinate ball relatively compact in $X\setminus A$, write $S=\ddc v$ there, and put $\eta=\ddc|z|^2$. Write $\nu=g\eta^n$, where $g$ is smooth and positive. Ordinary and non-pluripolar products agree on this ball, so $(\ddc v)^n=mg\eta^n$. The decreasing radial regularizations $v_\varepsilon$ from \cref{prop:bounded-monge-ampere-linear-density} satisfy
\[
 (\ddc v_\varepsilon)^n=F_\varepsilon\eta^n,
 \qquad
 F_\varepsilon\geq
       \bigl((mg)^{1/n}*\rho_\varepsilon\bigr)^n.
\]
The lower bound converges uniformly on smaller balls to $mg$. For a nonnegative smooth compactly supported $\chi$ there, it follows that
\begin{equation}\label{eq:mixed-trace-local-root-limit}
 \liminf_{\varepsilon\downarrow0}
       \int\chi\sqrt{F_\varepsilon g}\,\eta^n
                  \geq\sqrt m\int\chi\nu.
\end{equation}
The two trace integrals against $\chi$ converge by ordinary current convergence and mixed bounded Bedford--Taylor continuity. No identification of the absolutely continuous Hessian with the Monge--Amp\`ere density is needed.

For a smooth compactly supported $0\leq\chi\leq1$ on $X\setminus A$, choose a finite nonnegative partition $\chi=\sum_l\chi_l$ subordinate to such balls, and regularize separately in them. Write $a_{l,\varepsilon},b_{l,\varepsilon}$ for the integrals of the two smooth traces against $\chi_l$. The pointwise estimate and Cauchy--Schwarz, first on each ball and then on the finite sum, give
\[
 \left(\sum_l a_{l,\varepsilon}\right)
 \left(\sum_l b_{l,\varepsilon}\right)
 \geq\frac1n\left(\sum_l
       \int\chi_l\sqrt{F_{l,\varepsilon}g_l}\,\eta_l^n\right)^2.
\]
Take limits using \eqref{eq:mixed-trace-local-root-limit}. The two limiting trace integrals are bounded above by the corresponding full integrals in \eqref{eq:prescribed-density-mixed-trace}. Exhaust $X\setminus A$ with such cutoffs tending to one. The set $A$ has zero smooth volume, hence zero $\nu$ mass. This proves the estimate. Global non-pluripolar finiteness makes its second factor finite; the first is the pairing of a current with a smooth form on compact $X$. For $n=1$ the same argument uses the empty product, or follows directly from $S\geq\langle S\rangle=m\nu$.
\end{proof}
\index{Monge--Amp\`ere equation!mixed trace estimate}

\Needspace{12\baselineskip}
\begin{theorem}\label{thm:positive-nonpluripolar-mass-bigness}
If a closed positive $(1,1)$-current $T$ satisfies $\int_X\langle T^n\rangle>0$, its class is big. More generally, a pseudo-effective class is big if and only if its volume is positive.
\end{theorem}
\begin{proof}
Write $\alpha=\{\theta\}$ with $\theta$ smooth and closed. It suffices to prove bigness when $\operatorname{vol}(\alpha)>0$: the mass bound in \cref{prop:positive-mixed-mass-bound} gives this positivity from the asserted current, while that proposition already proves positive volume for big classes. Set $m=\operatorname{vol}(\alpha)>0$. Choose $C>0$ so that $\theta+C\omega>0$, and put
\[
 C_0=\int_X(\theta+(C+1)\omega)^{n-1}\wedge\omega>0.
\]

For $0<\varepsilon\leq1$ let $\theta_\varepsilon=\theta+\varepsilon\omega$. Pseudo-effectivity supplies $q\in\PSH(X,\theta)$; its current plus $\varepsilon\omega$ is strictly positive, so $\{\theta_\varepsilon\}$ is big by definition. Its normalized minimal potential $V_{\theta_\varepsilon}$ is a model potential: the largest nonpositive potential is its own singularity envelope. Its mass
\[
 m_\varepsilon=\operatorname{vol}(\alpha+\varepsilon\{\omega\})
                                                        \geq m
\]
is positive by class monotonicity and the decreasing boundary definition of volume.

Fix any smooth strictly positive $\ddc$-closed $(n-1,n-1)$-form $\tau$, normalized by $\int_X\omega\wedge\tau=1$. The measure $m_\varepsilon\omega\wedge\tau$ has a smooth density relative to $\omega^n$. Apply \cref{cor:relative-lp-capacity-boundedness} with reference $V_{\theta_\varepsilon}$ to obtain $u_\varepsilon$ such that
\[
 \begin{gathered}
 S_\varepsilon=\theta_\varepsilon+\ddc u_\varepsilon\geq0,
 \qquad \langle S_\varepsilon^n\rangle
                =m_\varepsilon\omega\wedge\tau,\\
 V_{\theta_\varepsilon}-B_\varepsilon
                   \leq u_\varepsilon\leq V_{\theta_\varepsilon}.
 \end{gathered}
\]
Here $B_\varepsilon$ is finite; no uniform bound for it is required. The gentle analytic strict current supplied by \cref{lem:logarithmic-cutoff-sum} in this big class has a proper analytic pole set $A_\varepsilon$. Its potential is a lower bound for $V_{\theta_\varepsilon}$ up to a constant, so $u_\varepsilon$ is locally bounded off $A_\varepsilon$.

The preceding lemma and \cref{thm:compact-nonpluripolar-finiteness} now give
\[
 \begin{aligned}
 \left(\int_X\theta\wedge\tau+\varepsilon\right)C_0
 &\geq\left(\int_X S_\varepsilon\wedge\tau\right)
        \left(\int_X\langle S_\varepsilon^{n-1}\rangle\wedge\omega\right)
 \\&\geq\frac{m_\varepsilon}{n}\geq\frac mn.
 \end{aligned}
\]
Indeed $\ddc\tau=0$ identifies the first trace with the displayed cohomological pairing. For the second trace use the shifted positive background $\theta_\varepsilon+C\omega\leq\theta+(C+1)\omega$ in the smooth mass bound; its positive powers have integrals at most $C_0$, uniformly in $\varepsilon$ and $\tau$. When $n=1$ the second factor is exactly $C_0=\int_X\omega$. Let $\varepsilon\downarrow0$ and undo the normalization of $\tau$. We obtain \eqref{eq:lamari-strict-tests} with $\delta=m/(nC_0)>0$. Lamari's criterion produces a current in $\alpha$ at least $\delta\omega$, proving bigness.
\end{proof}
\index{big class!positive non-pluripolar mass}\index{volume!bigness criterion}

The zero-mass alternative extends the obstacle concentration theorem to the boundary of the pseudo-effective cone.

\begin{corollary}\label{cor:pseudoeffective-obstacle-contact}
Let $\theta$ be a smooth closed real $(1,1)$-form and let $f:X\to[-\infty,+\infty]$ be quasi-continuous. Suppose its obstacle family $\mathcal A_\theta(f)$ is nonempty and uniformly bounded above. Then $u=P_\theta(f)$ satisfies
\[
 \mu_u\bigl(X\setminus\{u=f\}\bigr)=0.
\]
If $\{\theta\}$ is not big, $\mu_u=0$. For $r\geq1$, in every pseudo-effective class a nonempty rooftop $w=P_\theta(\min_{1\leq j\leq r}u_j)$ of $\theta$-psh potentials satisfies
\[
 \mu_w\leq\sum_{j=1}^r\mathbf1_{\{w=u_j\}}\mu_{u_j}.
\]
\end{corollary}
\begin{proof}
The envelope is $\theta$-psh by \cref{prop:obstacle-envelope-maximal-member}. In a non-big class the preceding theorem forbids positive total mass for any such potential, so its positive measure is zero. Suppose the class is big. Choose $C>\sup_Xu$ and put $g=\min\{f,C\}$. Since $u\leq f$ quasi-everywhere and $u<C$ everywhere, $u$ is a competitor for $g$. Conversely every competitor for $g$ is one for $f$, so $P_\theta(g)=u$. The function $g$ is quasi-continuous and bounded above; \cref{thm:quasicontinuous-obstacle-contact} applies. Its contact set $\{u=g\}$ equals $\{u=f\}$, because $u$ never reaches $C$. This proves concentration even when $f$ is unbounded above.

The finite minimum of the $u_j$ is quasi-continuous and bounded above, and its envelope is the specified rooftop. Concentration puts $\mu_w$ on the union of the finite contact sets. On each of them \cref{lem:ordered-potential-contact-measure} bounds its restriction by $\mu_{u_j}$. Adding those bounds proves the rooftop inequality; common poles have zero measure.
\end{proof}
\index{envelope!pseudo-effective contact concentration}
